% !TEX TS-program = pdflatexmk
\documentclass[12pt]{amsart}

%\usepackage[parfill]{parskip}    % Activate to begin paragraphs with an empty line rather than an indent

\usepackage[margin=1in]{geometry}

\usepackage{amsmath,amssymb,amsthm,latexsym,graphicx}
\usepackage[normalem]{ulem}
\usepackage{setspace} %used for doublespacing, etc.
\usepackage{hyperref}
\usepackage[shortlabels]{enumitem}
\usepackage{cancel}
\usepackage[dvipsnames,usenames]{color}
\usepackage[all]{xy}
\DeclareGraphicsRule{.tif}{png}{.png}{`convert #1 `dirname #1`/`basename #1 .tif`.png}

%Some useful environments.
\newtheorem{theorem}{Theorem}
\newtheorem{corollary}[theorem]{Corollary}
\newtheorem{conjecture}[theorem]{Conjecture}
\newtheorem{lemma}[theorem]{Lemma}
\newtheorem{proposition}[theorem]{Proposition}
\newtheorem{definition}[theorem]{Definition}
\newtheorem{example}[theorem]{Example}
\newtheorem{axiom}{Axiom}
\theoremstyle{remark}
\newtheorem{remark}{Remark}
\newtheorem*{exercise}{Exercise}%[section]

%Some useful shortcuts for our favorite fields
\newcommand{\RR}{{\mathbb R}}
\newcommand{\NN}{{\mathbb N}}
\newcommand{\ZZ}{{\mathbb Z}}
\newcommand{\QQ}{{\mathbb Q}}
\newcommand{\CC}{{\mathbb C}}

%Some useful shortcuts for formatting lists
\newcommand{\bc}{\begin{center}}
\newcommand{\ec}{\end{center}}
\newcommand{\be}{\begin{enumerate}}
\newcommand{\ee}{\end{enumerate}}
\newcommand{\bi}{\begin{itemize}}
\newcommand{\ei}{\end{itemize}}

%Some useful shortcuts for formatting mathematical symbols
\newcommand{\ol}[1]{\overline{#1}}
\newcommand{\oimp}[1]{\overset{#1}{\Longleftrightarrow}} %labeled iff symbol
\newcommand{\bv}[1]{\ensuremath{ \vec{\mathbf{#1}}} } %makes a vector.
\newcommand{\mc}[1]{\ensuremath{\mathcal{#1}}} %put something in caligraphic font
\newcommand{\mpg}[1]{\marginpar{ #1}} %to put comments in margins
\newcommand{\bsl}[1]{\texttt{\symbol{92}{\em #1}}} %for backslashes.
\newcommand{\normale}{\trianglelefteq}
\newcommand{\normal}{\triangleleft}

%Commenting tools
\usepackage{soul}
\definecolor{highlight}{rgb}{1,0.6,0.6}
\sethlcolor{highlight}
\newcommand{\hlm}[1]{\colorbox{highlight}{$\displaystyle #1$}}
\newtheoremstyle{mycomment}{\topsep}{-0in}{\small \itshape \sffamily}{}{\small \itshape\sffamily}{:}{.5em}{}
\theoremstyle{mycomment}
\newtheorem*{acomment}{\color{BrickRed}{Comment}}
\newcommand{\com}[1]{{\color{OliveGreen}\begin{acomment}{#1} %#2 \color{black} 
\end{acomment}\noindent}}
%\newcommand{\com}[1]{{\color{BrickRed}{\\ Comment:}\color{OliveGreen}{#1} \\}}
\newcommand{\red}[1]{{\color{BrickRed} #1}}
\def\midd{\,\middle\vert\, }
\newcommand{\blue}[1]{{\color{MidnightBlue}#1}}
\newcommand{\green}[1]{{\color{OliveGreen}#1}}
\newcommand{\mwrong}[2]{\red{\cancel{#1}}\green{#2}}
\newcommand{\wrong}[2]{\red{\sout{#1}}\green{#2}}
\definecolor{OliveGreen}{rgb}{.3,.5,.2}
\definecolor{MidnightBlue}{rgb}{.3,.4,.6}

\title{Assignment \#8 10/28/22}

\author{Coordinator: Stefano Fochesatto}
\begin{document}
\maketitle
%\setstretch{2.5} %Use for 2.5 spacing
\doublespacing %Use for double spacing

\section*{Section 5.1}
\begin{exercise}[\textbf{5.1.4} --- Stefano Fochesatto] Let $A$ and $B$ be finite groups and let $p$ be prime. Prove that any Sylow p-subgroup of $A\times B$
  is of the form $P\times Q$, where $P \in Syl_p(A)$ and $Q \in Syl_p(B)$. Prove that $n_p(A\times B) = n_p(A)n_p(B)$. Generalize both of these results 
  to a direct product of any finite number of finite groups.
  \begin{proof} Suppose $A$ and $B$ are finite groups and let $P$ be prime. We know that $|A| = p^{\alpha}m$ and $|B| = p^{\beta}n$ where $p \nmid m, n$.
    Note that by properties of direct products we know that $|A \times B| = p^{\alpha + \beta}mn$. Consider $P \in Syl_p(A)$ and $Q \in Syl_p(B)$ and note that $|P \times Q| = p^{\alpha + \beta}mn$ and $p \nmid mn$, therefore $P \times Q \in Syl_p(A \times B)$. By Sylow's Theorem we know every $H \in Syl_p(A\times B)$ is conjugate to $P \times Q$. Let $g \in A$ and $h \in B$ then there exists some $(g, h) \in (A \times B)$ such that $H = (g, h)(P, Q)(g, h)^{-1} = (gPg^{-1}, hQh^{-1})$. By Sylow's Theorem, we also know that $gPg^{-1} \in Syl_p(A)$ and similarly $hQh^{-1} \in Syl_p(B)$. Therefore every $H \in Syl_p(A\times B)$ is of the form $(P\times Q)$ where $P \in Syl_p(A)$ and $Q \in Syl_p(B)$.


    Since every $H \in Syl_p(A\times B)$ is of the form $(P\times Q)$ where $P \in Syl_p(A)$ and $Q \in Syl_p(B)$ we know that there are only 
    $n_p(A)n_p(B)$ ways to combine Sylow p-subgroups to get a $(P\times Q)$ product. Thus $n_p(AxB) = n_p(A)n_p(B)$.


    Let $G_1, G_2, \dots G_n$ be groups and let $G = G_1 \times  G_2 \times  \dots \times  G_n$, then any $X \in Syl_p(G)$ is of the form 
    $X = P_1 \times  P_2 \times \dots \times P_n$ where $P_i \in Syl_p(G_i)$ and $n_p(G) = \prod n_p(G_i)$. 

  \end{proof}
\end{exercise}
\vspace{.15in}



\begin{exercise}[\textbf{5.1.5} --- Glen Woodworth] Exhibit a nonnormal subgroup of $Q_8\times Z_4$ (note that every subgroup of each factor is normal).
  \begin{proof} Consider the subgroup $\langle (-k,0)\rangle=\{(-k,0),(1,0)\}$ in $Q_8\times Z_4$. Observe, \begin{align*}
  (j,0)(-k,0)(-j,0)&=(j(-k),0)(-j,0)=(-jk,0)(-j,0)=(i,0)(-j,0)\\
  &=(i(-j),0)=(-ij,0)=(k,0)\notin\langle (-k,0)\rangle.\end{align*}
 Hence, $\langle (-k,0)\rangle$ is not a normal subgroup of $Q_8\times Z_4$.
  \end{proof}
\end{exercise}
\vspace{.15in}

\begin{exercise}[\textbf{5.1.10} --- Noah Palmer] Let $p$ be a prime. Let $A$ and $B$ be two cyclic groups of order $p$ with generators $x$ and $y$, respectively. Set $E=A\times B$ so that $E$ is the elementary abelian group of order $p^2: E_{p^2}$. Prove that the distinct subgroups of $E$ of order $p$ are
  $$\langle x\rangle,\hspace{0.5cm}\langle xy\rangle,\hspace{0.5cm}\ldots \hspace{0.5cm},\langle xy^{p-1}\rangle,\hspace{0.5cm}\langle y\rangle$$
  \begin{proof}
  Let $p$ be a prime, $\langle x\rangle=A$ and $\langle y\rangle=B$ be two cyclic groups of order $p$, and set $E=A\times B$ so that $E$ is the elementary abelian group of order $p^2$. Since $A$ and $B$ are cyclic groups of prime order it follows that each element of $A$ and $B$ have order $p$. Hence, each of the subgroups
  $$\langle x\rangle,\hspace{0.5cm}\langle xy\rangle,\hspace{0.5cm}\ldots \hspace{0.5cm},\langle xy^{p-1}\rangle,\hspace{0.5cm}\langle y\rangle$$
  will have order $p$. Additionally, each of these groups are unique since each $y^{i}$ is unique within its own cyclic group which implies that the subgroups generated by $\langle(1,1,\ldots,x,\ldots, y^i,\ldots 1)\rangle$ will be unique for each $1\leq i<p$. From example $3$ in section $5.1$ of Dummit and Foote the number of subgroups of order $p$ in an elementary abelian group of order $p^2$ is $p+1$. Since, we have given a list of exactly $p+1$ distinct groups, we're done. 
  \end{proof}
  \end{exercise}
\vspace{.15in}





\section*{Section 5.2}
\begin{exercise}[\textbf{5.2.2} --- Oleksandr Bobrovnikov]
  In each of parts (a) to (d) give the lists of invariant factors for all abelian groups of the specified order:
  \begin{enumerate}[(a)]
    \item order 270:
    \begin{proof}[Solution:]
      Note that $270=2\times 3^3\times 5$. Since $2\times 3\times 5\mid n_1$, then $n_1\in\left\{ 2\times 3\times 5,2\times 3^2\times 5,2\times 3^3\times 5 \right\}$.
      If $n_1=2\times 3\times 5$ then from $n_{i+1}\mid n_i$ we get $n_2=3$ and $n_3=3$.
      If $n_1=2\times 3^2\times 5$ then similarly we get $n_2=3$.
      If $n_1=2\times 3^3\times 5$ then this is the only invariant factor.
    \end{proof}
    \item order 9801:
    \begin{proof}[Solution:]
      Note that $9801=3^4\times 11^2$. Thus\\
      $n_1\in\left\{ 3\times 11,3^2\times 11,3^3\times 11,3^4\times 11,3\times 11^2,3^2\times 11^2,3^3\times 11^2,3^4\times 11^2\right\}$.
  
      If $n_1=3\times 11$ then $n_2=3\times 11$ since otherwise we will not obtain $n=\prod n_j$ (if we take 3 for example, there will be not enough 11).
      Then list of invariant factors is $3\times 11,3\times 11,3,3$.
      
      If $n_1=3^2\times 11$ then similarly $n_2$ should be $3\times 11$ or $3^2\times 11$. Then possible lists are the following: $3^2\times 11,3\times 11,3$ and $3^2\times 11,3^2\times 11$.
  
      If $n_1=3^3\times 11$ then the only possible list is $3^3\times 11,3\times 11$
  
      If $n_1=3^4\times 11$ then $n_2=11$.
  
      If $n_1=3\times 11^2$ then the only possible list is $n_1,3,3,3$.
  
      If $n_1=3^2\times 11^2$ then possible lists are $n_1,3,3$ and $n_1,3^2$.
  
      If $n_1=3^4\times 11^2$ then the list is $n_1,3$.
  
      And finally if $n_1=9801$, then this is the entire list.
    \end{proof}
    \item order 320:
    \begin{proof}[Solution:]
      Note that $320 = 2^6\times 5$. So $n_1=2^j\times 5$, where $1\leq j\leq 6$. So possible list are:
      $\left( 2\times 5,2,2,2,2,2 \right)$, $\left( 2^2\times 5, 2^2,2^2 \right)$, $\left( 2^2\times 5,2^2,2,2 \right)$,
      $(2^2\times 5,2,2,2,2)$, $(2^3\times 5,2^3)$, $(2^3\times 5, 2^2,2)$, $(2^3\times 5, 2,2,2)$, $(2^4\times 5,2^2)$, $(2^4\times 5,2,2)$, $(2^5\times 5,5)$, $(2^6)$
    \end{proof}
    \item order 105:
    \begin{proof}[Solution:]
      Note that $105=3\times 5\times 7$. Then the only possible list is $n_1=105$.
    \end{proof}
  \end{enumerate}
\end{exercise}
\vspace{.15in}


\begin{exercise}[\textbf{5.2.4} --- Noah Palmer] In each of parts $(a)$ to $(d)$ determine which pairs of abelian groups listed are isomorphic.
  \begin{enumerate}[(a)]
  \item $\{4,9\}$, $\{6,6\}$, $\{8,3\}$, $\{9,4\}$, $\{6,4\}$, $\{64\}$.
  \item $\{2^2, 2\cdot3^2\}$, $\{2^2\cdot3, 2\cdot3\}$, $\{2^3\cdot3^2\}$, $\{2^2\cdot3^2,2\}$.
  \item $\{5^2\cdot7^2,3^2\cdot5\cdot7\}$, $\{3^2\cdot5^2\cdot7,5\cdot7^2\}$, $\{3\cdot5^2,7^2,3\cdot5\cdot7\}$, $\{5^2\cdot7,3^2\cdot5,7^2\}$.
  \item $\{2^2\cdot5\cdot7, 2^3\cdot5^3,2\cdot5^2\}$, $\{2^3\cdot5^3\cdot7,2^3\cdot5^3\}$, $\{2^2,2\cdot7,2^3,5^3,5^3\}$, $\{2\cdot5^3,2^2\cdot5^3,2^3,7\}$.
  \end{enumerate}
  \begin{proof}
  We use the fact that $Z_m\times Z_n\cong Z_{mn}$ if and only if $\gcd(m,n)=1$ to rearrange the abelian groups into isomorphic pairs, where possible.
  \begin{enumerate}[(a)]
  \item Note that $ Z_4\times Z_9\cong Z_9\times Z_4$. Thus, $\{4,9\}\cong\{9,4\}$.
  \item Note that $ Z_{2^2}\times Z_{2\cdot3^2}\cong Z_{2^2\cdot3^2}\times Z_2$. Thus, $\{2^2, 2\cdot3^2\}\cong\{2^2\cdot3^2,2\}$.
  \item Note that $Z_{5^2\cdot 7^2}\times Z_{3^2\cdot 5\cdot 7}\cong Z_{3^2\cdot 5^2\cdot 7}\times Z_{5\cdot 7^2}$. Also, $Z_{3^2\cdot 5^2\cdot 7}\times Z_{5\cdot 7^2}\cong Z_{5^2\cdot 7}\times Z_{3^2\cdot 5}\times Z_{7^2}$. Thus, $\{5^2\cdot7^2,3^2\cdot5\cdot7\}\cong\{3^2\cdot5^2\cdot7,5\cdot7^2\}\cong\{5^2\cdot7,3^2\cdot5,7^2\}$.
  \item Note that $ Z_{2^2}\times Z_{2\cdot7}\times Z_{2^3}\times Z_{5^3}\times Z_{5^3}\cong Z_{2^2\cdot5^3}\times Z_{2\cdot7\cdot5^3}\times Z_{2^3}\cong Z_{2\cdot5^3}\times Z_{2^2\cdot5^3}\times Z_{2^3}\times Z_7$. Therefore, $\{2^2,2\cdot7,2^3,5^3,5^3\}\cong\{2\cdot5^3,2^2\cdot5^3,2^3,7\}$.
  \end{enumerate}
  \end{proof}
  \end{exercise}
\vspace{.15in}

  


  \begin{exercise}[\textbf{5.2.9} --- Noah Palmer] Let $A=Z_{60}\times Z_{45}\times Z_{12}\times Z_{36}$. Find the number of elements of order $2$ and the number of subgroups of index $2$ in $A$.
    \begin{proof}
    Begin by rewriting the direct product in its invariant form $Z_{60}\times Z_{45}\times Z_{12}\times Z_{36}\cong Z_{180}\times Z_{60}\times Z_3\times Z_{36}\cong Z_{180}\times Z_{180}\times Z_{12}\times Z_{3}$. Now, note that each of $Z_{180}$ and $Z_{12}$ has an element of order $2$ whereas $Z_{3}$ does not. Hence, there will be $2^3$ elements of order less than or equal to $2$. From these we must exclude the identity since it has order $1$. Thus, there are $7$ elements of order $2$.
    
    To find the subgroups of index $2$ in $A$ consider the elementary divisor factorization of $A=Z_{5}\times Z_{5}\times Z_{4}\times Z_{4}\times Z_4\times Z_{9}\times Z_{9}\times Z_3\times Z_3$. Now, any subgroup of order $583,200$ will have index $2$ in $A$. From the above factorization, and the prime factorization of $583,200$, these subgroups will have the form $A^\prime\times Z_{5}\times Z_{5}\times Z_{9}\times Z_{9}\times Z_3\times Z_3$ where $A^\prime$ is a subgroup of $Z_{4}\times Z_{4}\times Z_4$ with order $32$. There are three ways to select such a subgroup (e.g. $Z_4\times Z_4\times Z_2$, $Z_2\times Z_4\times Z_4$, and $Z_4\times Z_2\times Z_4$). Thus, there are three subgroups of index $2$ in $A$.
    \begin{center}
    
    \end{center}
    \end{proof}
    \end{exercise}
\vspace{.15in}

  




\section*{Section 5.4}
\begin{exercise}[\textbf{5.4.1} --- Stefano Fochesatto] Prove that if $x, y \in G$ then $[y, x]= [x, y]^{-1}$. Deduce that for any subsets $A, B$ of $G$, $[A, B] = [B, A]$
  \begin{proof} Suppose if $x, y \in G$ then we know $[y, x] = y^{-1}x^{-1}yx$. Note that $y^{-1}x^{-1}yx = ((xy)^{-1})(yx) = ((yx)^{-1}(xy))^{-1} = (x^{-1}y^{-1}xy)^{-1} = [x, y]^{-1}$.

    Let $A, B \leq G$ and consider the group generated by $[A, B]$. By definition we know that
    $[A, B] = \langle [a, b] | a \in A, b \in B \rangle$. Since $[A, B]$ is a group generated by a set i.e the set of all commutators, it must also be generated by a set of commutator inverses, $[A, B] = \langle [a, b]^{-1} | a \in A, b \in B \rangle$
    By the previous result we know that $[a, b]^{-1} = [b, a]$ so $[A, B] = \langle [b, a] | a \in A, b \in B \rangle = [B, A]$. 
  \end{proof}
\end{exercise}
\vspace{.15in}





\begin{exercise}[\textbf{5.4.11} --- Oleksandr Bobrovnikov]
  Prove that if $G=HK$ where $H$ and $K$ are characteristic subgroups of $G$ with $H \cap K = 1$
  then $Aut(G) \cong Aut(H) \times Aut(K)$. Deduce that if $G$ is an abelian group of finite order then $Aut( G)$ is isomorphic to the direct product of the automorphism groups of its Sylow subgroups.
  \begin{proof}[Solution:]
    Since $H$ and $K$ are normal subgroups of $G$ with $H \cap K = 1$, we know that $HK\cong H\times K$.
    Let $G=H\times K$.
    We define $\varphi_H\in Aut(H)$  by $\varphi_H(h)=\pi_H(\varphi(h,1))$ and similarly we define $\varphi_K\in Aut(K)$.% (Here $\pi_H=\pi_{Aut(H)}$.)
    Then consider $\Phi:Aut(G)\to Aut(H) \times Aut(K)$, which acts by $\varphi\mapsto(\varphi_H,\varphi_K)$.
    This is an isomorphism.
    First note that it is a homomorphisms:
    \begin{multline}
      [\Psi(\varphi\circ \psi)](h,k)
      =\\((\varphi\circ\psi)_H(h),(\varphi\circ\psi)_K(k))
      =(\pi_H[(\varphi\circ\psi)(h,1)],
      \pi_K[(\varphi\circ\psi)(1,k)]
      )=\\
      ((\pi_H\circ\varphi\circ\pi_H\circ\psi)(h,1),
      (\pi_K\circ\varphi\circ\pi_K\circ\psi)(1,k))
      =
      ((\varphi_H\circ\psi_H)(h,1),
      (\varphi_K\circ\psi_K)(1,k))
      =\\
      [\Psi(\varphi)\circ\Psi(\psi)](h,k),
    \end{multline}
    and similarly $\Psi(\varphi^{-1})=\Psi(\varphi)^{-1}$.
    Finally, $\Psi$ is a bijection since by $\varphi_H$ and $\varphi_K$ one can uniquely determine $\varphi$.
    \\ If $G$ is an abelian group of finite order then every its Sylow $p$-subgroup are normal and thus are characteristic, and intersection of any two of them is identity.
    Then using induction by number of these subgroups one may show that 
    $Aut( G)$ is isomorphic to the direct product of the automorphism groups of its Sylow subgroups.
    Base is proven above, and a step uses base and associativity of direct product. 


  \end{proof}
\end{exercise}
\vspace{.15in}




\begin{exercise}[\textbf{5.4.17} --- Glen Woodworth] If $K$ is normal subgroup of $G$ and $K$ is cyclic, prove that $G'\le C_G(K)$. [Recall that the automorphism group of a cyclic group is abelian.]
  \begin{proof} Let $K\trianglelefteq G$. Recall that the automorphism group of a cyclic group is abelian. Since $K$ is cyclic, Aut$(K)$ is abelian. Consider the map
  \begin{align*}
  \Phi:G&\rightarrow \text{Aut }(K)\\
  g&\mapsto\phi_g:k\mapsto gkg^{-1}.\end{align*}
  Since $K\trianglelefteq G$, $gkg^{-1}\in K$ for all $k\in K$ and for all $g\in G$. Thus, $\Phi$ is well-defined. We now show that $\Phi$ is a homomorphism. Let $g_1,g_2\in G$ and $k\in K$. Observe,
  \begin{align*}
  \Phi(g_1g_2)(k)&=\phi_{g_1g_2}(k)=g_1g_2k(g_1g_2)^{-1}=g_1g_2kg_2^{-1}g_1^{-1}=\phi_{g_1}(g_2kg_2^{-1})\\
  &=\phi_{g_1}(\phi_{g_2}(k))=(\phi_{g_1}\circ\phi_{g_2})(k)=(\Phi(g_1)\circ\Phi(g_2))(k).
  \end{align*}
Thus, $\Phi$ is a homomorphism. We now show that $\ker\Phi= C_G(K)$. Let $x\in C_G(K)$. Then \[\Phi(x)(k)=xkx^{-1}=k.\]
Thus, $\Phi(x)=Id$. Suppose $y\in\ker\Phi$. We have \[k=\Phi(y)(k)=yky^{-1}.\] Therefore, $y\in C_G(K)$. Hence, $\ker\Phi=C_G(K)$. Proposition 7 from Section 5.4 states that if $\phi:G\rightarrow A$ is any homomorphism where $G$ is a group and $A$ is an abelian group, then $G'\le\ker\phi$. Since Aut$(K)$ is abelian, and $\Phi$ is a homomorphism with $\ker\Phi=C_G(K)$, we have $G'\le\ker\Phi=C_G(K)$.
  \end{proof}
\end{exercise}
\vspace{.15in}





\section*{Section 5.5}


\begin{exercise}[\textbf{5.5.1} --- Oleksandr Bobrovnikov]
  Prove that $C_K(H)=\ker\varphi$ (recall that $C_K(H)=C_G(H)\cap K$).
  \begin{proof}
    Note that $\ker \varphi=\left\{ (1,k)\in G\midd \forall h\in H:~k\cdot h =h \right\}$ and\\
    $C_K(H)=C_G(H)\cap K=\left\{ (1,k)\in G\midd\forall (h,s)\in G:~ (h,s)(1,k)=(1,k)(h,s)  \right\}$. The second is equivalent to
    $C_K(H)=\left\{ (1,k)\in G\midd\forall (h,s)\in G:~ (hs\cdot 1,sk)=(1k\cdot h,ks)  \right\}$.
    From this follows that $C_K(H)=\ker\varphi$.
  \end{proof}
\end{exercise}
\vspace{.15in}




\begin{exercise}[\textbf{5.5.2} --- Stefano Fochesatto] Let $H$ and $K$ be groups, let $\varphi$ be a homomorphism from $K$ into $Aut(H)$ and identify $H$ and $K$ as subgroups of $G = H\rtimes_{\varphi} K$. Prove that $C_H(K) = N_H(K)$
  \begin{proof} 
    By definition we know that $C_H(K) \subseteq N_H(K)$. Let $(h, 1) \in N_H(K)$. By definition we know that for all $(1, k_1) \in K$ there exists some $(1, k_2) \in K$ such that $(h,1)(1,k_1)(h^{-1},1) = (1,k_2)$. By multiplication in $G$ we get $(h,1)(\varphi_{k_1}(h^{-1}), k_1) = (1,k_2)$. Again multiplication in $G$ gives us, $(\varphi_{k_1}(h)\varphi_{k_1}(h^{-1}), k_1) = (1,k_2)$. So $k_1 = k_2$ and by we can get $(h, 1) \in C_H(K)$. Thus $C_H(K) = N_H(K)$.
  \end{proof}
\end{exercise}
\vspace{.15in}






\begin{exercise}[\textbf{5.5.8} --- Glen Woodworth] Construct a non-abelian group of order $75$. Classify all groups of order 75 (there are three of them). [Use Exercise 6 to show that the non-abelian group is unique.] (The classification of groups of order $pq^2$, where $p$ and $q$ are primes with $p<q$ and $p$ not dividing $q-1$, is quite similar.)
  \begin{proof} Consider the homomorphism \begin{align*}\Phi:Z_{25}&\rightarrow\text{Aut}(Z_{3})\\
\overline{0}&\mapsto\text{Id}:\overline{1}\mapsto\overline{1}\\
\overline{1}&\mapsto \phi:\overline{1}\mapsto \overline{2}.
  \end{align*} 
  We show that $\Phi$ is well-defined. Observe,\[\phi(\overline{1}+\overline{2})=\phi(\overline{0})=\overline{0}=\overline{2}+\overline{1}=\phi(\overline{1})+\phi(\overline{2})\]
Thus, $\phi$ is a homomorphism. Since $\phi$ is obviously a bijection, $\phi\in$Aut$(Z_3)$. Thus, $\Phi$ is well-defined. We now show $\Phi$ is a homomorphism. First, observe that since $\Phi(\overline{1})(\overline{1})=\phi(\overline{1})=\overline{2}$ and since there are only two options in the image, if $x\in Z_{25}$ is odd, then $\Phi(x)(\overline{1})=\phi(\overline{1})=\overline{2}$. Similarly, if $y\in Z_{25}$ is even, then $\Phi(y)(\overline{1})=$Id$(\overline{1})=\overline{1}$. Now, suppose $g_1,g_2\in Z_{25}$. We have three cases to consider. For the first case, suppose $g_1$ and $g_2$ are both even. Then $g_1+g_2$ is even. Observe,\[
\Phi(g_1+g_2)(\overline{1})=\text{Id}(\overline{1})=\text{Id}(\text{Id}(\overline{1}))=\Phi(g_1)(\Phi(g_2)(\overline{1}))=(\Phi(g_1)\circ\Phi(g_2))(\overline{1}).\]
For the second case, suppose $g_1$ and $g_2$ are both odd. Then $g_1+g_2$ is even. We have
\[\Phi(g_1+g_2)(\overline{1})=\text{Id}(\overline{1})=\overline{1}=\phi(\overline{2})=\phi(\phi(\overline{1}))=\Phi(g_1)(\Phi(g_2)(\overline{1}))=(\Phi(g_1)\circ\Phi(g_2))(\overline{1}).\]
For the last case, suppose $g_1$ is even and $g_2$ is odd. Then $g_1+g_2$ is odd. Observe,
\[\Phi(g_1+g_2)(\overline{1})=\phi(\overline{1})=\text{Id}(\phi(\overline{1}))=\Phi(g_1)(\Phi(g_2)(\overline{1}))=(\Phi(g_1)\circ\Phi(g_2))(\overline{1}).\]
Thus, $\Phi$ is a well-defined homomorphism. Hence, by the definition of semidirect product and Theorem 10, $G=Z_3\rtimes_\Phi Z_{25}$ is a group of order $75$. Let $\cdot$ denote the left action of $Z_{25}$ on $Z_3$ determined by $\Phi.$ To see that $G$ is nonabelian, consider $(\overline{1},\overline{3}),(\overline{2},\overline{4})\in G$. Observe,
\begin{align*}
(\overline{1},\overline{3})(\overline{2},\overline{4})(\overline{1},\overline{3})^{-1}&=(\overline{1}+\overline{3}\cdot\overline{2},\overline{3}+\overline{4})(\overline{22}\cdot\overline{2},\overline{22})=(\overline{1}+\Phi(\overline{3})(\overline{2}),\overline{7})(\Phi(\overline{22})(\overline{2}),\overline{22})\\
&=(\overline{1}+\phi(\overline{2}),\overline{7})(\text{Id}(\overline{2}),\overline{22})=(\overline{1}+\overline{1},\overline{7})(\overline{2},\overline{22})=(\overline{2},\overline{7})(\overline{2},\overline{22})\\
&=(\overline{2}+\overline{7}\cdot\overline{2},\overline{7}+\overline{22})=(\overline{2}+\Phi(\overline{7})(\overline{2}),\overline{4})=(\overline{2}+\phi(\overline{2}),\overline{4})=(\overline{2}+\overline{1},\overline{4})\\
&=(\overline{0},\overline{4}).
\end{align*}
Since $(\overline{2},\overline{4})\ne(\overline{0},\overline{4})$, $G$ is nonabelian. We are given that there are only three groups of order 75. The other two groups of order 75 are $Z_{75}$ and $Z_{15}\times Z_{5}$, which are both abelian. Note that since $(5,15)\ne 1$, $Z_{75}$ is not isomorphic to $Z_{15}\times Z_{5}$. Exercise 6 of section 5.5 states that if there exists another homomorphism $\Phi_1$ from $Z_{25}$ to Aut$(Z_3)$ such that $\Phi(Z_{25})$ and $\Phi_1(Z_{25})$ are conjugate subgroups of Aut$(Z_3)$, then $Z_3\rtimes_{\Phi}Z_{25}\cong Z_3\rtimes_{\Phi_1}Z_{25}.$ If we suppose that $\Phi_1$ is non-trivial, then $\phi\in\Phi_1(Z_{25})$ because $\phi$ is the only non-identity element in Aut$(Z_3)$. Since Id$\in\Phi_1(Z_{25})$ by the definition of a homomorphism, $\Phi_1(Z_{25})=\Phi(Z_{25})$. By Exercise 6, $Z_3\rtimes_{\Phi}Z_{25}\cong Z_3\rtimes_{\Phi_1}Z_{25}.$
  \end{proof}
\end{exercise}
\vspace{.15in}



\end{document}